\scanpage{97}\mspage{41}
% archive sheet number: 96
% L1
\unsure{15.11.|15,11.}

% L2-5
Je crois que j'ai enfin compris, cette nuit, la situation exacte. Il y a lieu, dans $\mathrm{Cat}_{/S}$ et dans les autres catégories des diagrammes p.~8 (s'envoyant toutes dans $\mathrm{Cat}_{/S}$)\unsure{,} d'envisager \deleted{trois} \inserted{quatre} sortes de ``$W$-équivalence''.

% L6-8
a) La $W$-\uline{équivalence} ``\uline{grossière}'' :
\diagram{\begin{tikzcd}[ampersand replacement=\&,column sep=tiny,row sep=tiny]
 f:X\arrow[rr]\arrow[dr,no head]\&\&Y\arrow[dl,no head]\\
 \&S\&
\end{tikzcd}}
\inserted{dans $\mathrm{Cat}_{/S}$} est une $W$-équivalence grossière si en \deleted{\unsure{tant}} que flèche de $\mathrm{Cat}$, c'est une $W$-équivalence.
\margin{Stable par chgt de base à la fois lisse et propre, plus \unsure{génér.} si $S'\longrightarrow S$ est une Hot-fibration}
% slanted

% L9-10
b) La $W$-\uline{équivalence} ``\uline{par fibres}\inserted{(sur $S$)}'' : $\forall s\in S$, $f_s:X_s\longrightarrow Y_s$ est une $W$-équivalence. \deleted{(Signifie que c'est une $W$-équiv, et le reste par \unsure{tt chgt}}
\margin{Stable par tt chgt de base $S'\longrightarrow S$}
% slanted

% L11-12
c) La \uline{$W$-équivalence} ``\deleted{\unsure{co}}\uline{locale}'' \uline{sur} $S$'' : $\forall s\in S$, $f_{/s}:X_{/s}\longrightarrow Y_{/s}$ est une $W$-équivalence. \inserted{Signifie que c'est une $W$-équiv., et le reste par \deleted{tout} \unsure{chgt} de base $S'\longrightarrow S$ qui est un isom. local.}
\margin{devrait être stable par chgt de base \uline{lisse} $S'\longrightarrow S$}
% slanted

% L13-14
\deleted{\unsure{d)}}\unsure{c')|d)} La $W$-\uline{équivalence} ``\uline{colocale sur} $S$'' : $\forall s\in S$, $ {}_{s\backslash}f:{}_{s\backslash}X\longrightarrow{}_{s\backslash}Y$ est une $W$-équivalence. \inserted{Signifie que c'est une $W$-équivalence et le reste par tt chgt. de base $S'\longrightarrow S$ qui est un co-isom local.}
\margin{devrait être stable \unsure{pour|par} chgt. de base \uline{propre} sur $S$}
% slanted

% L15-19
Je vais résumer dans un énoncé récapitulatif les relations entre ces quatre notions de $W$-équivalences, en pensant au cas $W=W_{\unsure{\infty|0}}$. Une fois la chose faite, je me préoccuperai de dégager les conditions minimales sur \inserted{un} $W$ \inserted{général} qui assurent que c'est vrai pour $W$.

% L20
\uline{Proposition 1} \unsure{.} Soit
\diagram{\begin{tikzcd}[ampersand replacement=\&,column sep=tiny,row sep=tiny]
 f:X\arrow[rr]\arrow[dr,no head]\&\&Y\arrow[dl,no head]\\
 \&S\&
\end{tikzcd}}
dans $\mathrm{Cat}_{/S}$.

% L21-23
1) La $W$-équivalence locale sur $S$, et la $W$-équivalence colocale sur $S$, impliquent la $W$-équivalence grossière :
\diagram{\begin{tikzcd}[ampersand replacement=\&,row sep=tiny,column sep=tiny]
 \deleted{\unsure{a}}c\arrow[dr,Rightarrow]\&\&c'\deleted{\unsure{d}}\arrow[dl,Rightarrow]\\
 \&a\&
\end{tikzcd}}

% L24-27
2) Si $X,Y$ sont propres sur $S$, ou $X,Y$ lisses sur $S$, ou $X,Y$ des Hot-fibrés sur $S$, alors la $W$-équivalence par fibres implique que c'est une $W$-équivalence \uline{universelle} (\inserted{sur $S$}) (\deleted{stable} $W$-équivalence, et le reste par \uline{tt} chg. de base $S'\longrightarrow S$) donc dans ces \unsure{trois} cas, on a
\margin{Il suffit $X,Y$ $W$-propres sur $S$, ou $W$-lisses sur $S$, ou $W$-fibrés sur $S$,}
% slanted
\diagram{\begin{tikzcd}[ampersand replacement=\&,column sep=tiny,row sep=tiny]
 \& b\arrow[dl,Rightarrow]\arrow[dr,Rightarrow]\&\\
 c\arrow[dr,Rightarrow]\&\&c'\arrow[dl,Rightarrow]\\
 \&a\&
\end{tikzcd}}
% L28-31
Mais on peut préciser \unsure{chacun} \deleted{\unsure{cas}} de ces trois cas :

\unsure{--}(2 fib) Si $X,Y$ sont Hot-fibrés sur $S$, alors la $W$-équivalence \inserted{gross.} de $f$ implique que c'est une $W$-équiv. par fibres, donc les quatre notions de $W$-équivalence sont équivalentes pour $f$.
\margin{$a\Rightarrow b$ seulement moyennant L~5, voir ci-dessous p.~43}
% slanted
\diagram{\begin{tikzcd}[ampersand replacement=\&,column sep=tiny]
 \unsure{a}\&\unsure{\Longleftarrow|\Longleftrightarrow}\&c\&\Longleftrightarrow\&c'\&\Longleftarrow\&b
\end{tikzcd}}
